% ===========================================================================
%  Programmation Frontend — Lab 6 : Consommer une API Django
%  ESP/UCAD — 2025–2026
%  Dr. El Hadji Bassirou TOURÉ
%  Compilé avec : xelatex (2–3 passes)
% ===========================================================================
\documentclass[11pt,a4paper]{article}

\usepackage{fontspec}
\setmainfont{Latin Modern Roman}
\setsansfont{Latin Modern Sans}
\setmonofont{Latin Modern Mono}
\usepackage{polyglossia}
\setdefaultlanguage{french}

\usepackage[margin=2.5cm]{geometry}
\usepackage{amsmath,amssymb}
\usepackage{graphicx}
\usepackage{xcolor}
\usepackage{booktabs}
\usepackage{hyperref}
\usepackage{tcolorbox}
\tcbuselibrary{theorems,breakable,skins}
\usepackage{enumitem}
\usepackage{fancyhdr}
\usepackage{array}
\usepackage{pifont}

\newcommand{\cmark}{\ding{51}}
\newcommand{\xmark}{\ding{55}}

\hypersetup{colorlinks=true,linkcolor=blue!60!black,urlcolor=blue!60!black,citecolor=blue!60!black}

\newtcolorbox{objectifbox}[1][]{colback=blue!5,colframe=blue!50,fonttitle=\bfseries,title=#1,breakable}
\newtcolorbox{infobox}[1][]{colback=green!5,colframe=green!50!black,fonttitle=\bfseries,title=#1,breakable}
\newtcolorbox{attentionbox}[1][]{colback=orange!5,colframe=orange!70,fonttitle=\bfseries,title=#1,breakable}
\newtcolorbox{retenirbox}[1][]{colback=yellow!10,colframe=yellow!50!black,fonttitle=\bfseries,title=#1,breakable}
\newtcolorbox{prereqbox}[1][]{colback=cyan!5,colframe=cyan!50!black,fonttitle=\bfseries,title=#1,breakable}
\newtcolorbox{projetbox}[1][]{colback=teal!5,colframe=teal!50!black,fonttitle=\bfseries,title=#1,breakable}

\usepackage{listings}
\lstdefinelanguage{JavaScript}{keywords={break,case,catch,continue,debugger,default,delete,do,else,finally,for,function,if,in,instanceof,new,return,switch,this,throw,try,typeof,var,void,while,with,const,let,class,export,extends,import,super,yield,async,await,of,from,true,false,null,undefined},morecomment=[l]{//},morecomment=[s]{/*}{*/},morestring=[b]',morestring=[b]",morestring=[b]`,sensitive=true}
\lstset{basicstyle=\ttfamily\small,backgroundcolor=\color{gray!8},frame=single,framerule=0.5pt,rulecolor=\color{gray!40},breaklines=true,breakatwhitespace=true,showstringspaces=false,tabsize=2,xleftmargin=4pt,xrightmargin=4pt,aboveskip=8pt,belowskip=8pt,language=JavaScript,keywordstyle=\color{blue!70!black}\bfseries,commentstyle=\color{green!50!black}\itshape,stringstyle=\color{red!60!black},literate={é}{{\'e}}1 {è}{{\`e}}1 {ê}{{\^e}}1 {à}{{\`a}}1 {ù}{{\`u}}1 {ô}{{\^o}}1 {î}{{\^i}}1 {ç}{{\c{c}}}1 {É}{{\'E}}1 {→}{{$\rightarrow$}}1}
\lstdefinestyle{terminal}{language={},keywordstyle={},commentstyle={},stringstyle={},basicstyle=\ttfamily\small,backgroundcolor=\color{black!5}}

\pagestyle{fancy}\fancyhf{}
\fancyhead[L]{\small Prog. Frontend --- Lab 6 : API Django}
\fancyhead[R]{\small ESP/UCAD}
\fancyfoot[C]{\thepage}
\fancyfoot[L]{\small Dr. El Hadji Bassirou TOURÉ}
\fancyfoot[R]{\small 2025--2026}
\renewcommand{\headrulewidth}{0.4pt}
\renewcommand{\footrulewidth}{0.4pt}

\title{\vspace{-1cm}\textbf{\LARGE Programmation Frontend}\\[0.8cm]\textbf{\huge Lab 6 --- Consommer une API Django}\\[0.3cm]\textbf{\Large CORS --- authentification --- CRUD complet}\\[0.5cm]\large ESP/UCAD --- 2025--2026}
\author{Dr. El Hadji Bassirou TOURÉ\\Département Génie Informatique\\Université Cheikh Anta Diop de Dakar}
\date{}

\begin{document}
\maketitle
\thispagestyle{empty}
\vspace{0.5cm}

\begin{prereqbox}[Informations pratiques]
\begin{center}
\begin{tabular}{lp{10cm}}
\toprule
\textbf{Durée estimée} & 2h (une séance) \\
\textbf{Prérequis} & Avoir terminé le Lab 5 (Next.js App Router). Python 3.10+ installé. \\
\textbf{Fichiers de départ} & Votre projet \texttt{annuaire-next/} du Lab 5 + le script \texttt{setup\_django.sh} fourni \\
\textbf{Livrables} & Annuaire Next.js connecté au backend Django avec login et CRUD \\
\textbf{Validation} & Les données sont persistées : ajouter un étudiant, recharger, il est toujours là. \\
\bottomrule
\end{tabular}
\end{center}
\end{prereqbox}

\begin{objectifbox}[Objectifs du lab]
À l'issue de ce lab, chaque étudiant sera capable de :
\begin{enumerate}[nosep]
\item Expliquer ce qu'est \textbf{CORS} et pourquoi le navigateur bloque certaines requêtes
\item Utiliser des \textbf{variables d'environnement} pour configurer l'URL de l'API
\item Centraliser les appels API dans un module \texttt{lib/api.js}
\item Implémenter l'\textbf{authentification par token} (login + envoi dans les headers)
\item Réaliser un \textbf{CRUD complet} (Create, Read, Delete) avec des données persistées
\end{enumerate}
\end{objectifbox}

\tableofcontents
\newpage

% ===========================================================================
\section{Partie 1 --- Installer et explorer le backend}
% ===========================================================================

\begin{objectifbox}[Objectif --- 10 minutes]
Lancer le backend Django fourni et tester l'API avec \texttt{curl}.
\end{objectifbox}

\subsection{Lancer le backend}

L'enseignant vous fournit un script \texttt{setup\_django.sh}. Ce script crée un projet Django complet avec une base de données pré-remplie de 10 étudiants sénégalais.

\begin{lstlisting}[style=terminal]
chmod +x setup_django.sh
./setup_django.sh
\end{lstlisting}

Puis lancez le serveur Django :

\begin{lstlisting}[style=terminal]
cd annuaire-api
source venv/bin/activate
python manage.py runserver
\end{lstlisting}

Le serveur tourne sur \texttt{http://localhost:8000}. \textbf{Gardez ce terminal ouvert} pendant tout le lab.

\subsection{Les endpoints de l'API}

\begin{center}
\begin{tabular}{llp{4.5cm}l}
\toprule
\textbf{Méthode} & \textbf{URL} & \textbf{Description} & \textbf{Auth} \\
\midrule
POST & \texttt{/api/auth/login/} & Connexion $\to$ token & Non \\
GET & \texttt{/api/students/} & Liste des étudiants & Token \\
POST & \texttt{/api/students/} & Créer un étudiant & Token \\
GET & \texttt{/api/students/3/} & Détail étudiant 3 & Token \\
DELETE & \texttt{/api/students/3/} & Supprimer étudiant 3 & Token \\
\bottomrule
\end{tabular}
\end{center}

\subsection{Tester avec \texttt{curl}}

Ouvrez un \textbf{deuxième terminal}. Commencez par le login :

\begin{lstlisting}[style=terminal]
curl -X POST http://localhost:8000/api/auth/login/ \
  -H "Content-Type: application/json" \
  -d '{"username":"prof","password":"le mot de passe défini dans setup_django.sh"}'
\end{lstlisting}

Résultat :

\begin{lstlisting}[style=terminal]
{"token":"abc123def456...","username":"prof"}
\end{lstlisting}

\textbf{Copiez le token.} Utilisez-le pour accéder aux données :

\begin{lstlisting}[style=terminal]
curl http://localhost:8000/api/students/ \
  -H "Authorization: Token abc123def456..."
\end{lstlisting}

Vous devez voir la liste des 10 étudiants au format JSON : Fatou Ndiaye, Moussa Diop, Aminata Sow...

\begin{attentionbox}[Sans le token --- erreur 401]
Si vous oubliez le header \texttt{Authorization}, le backend répond :

\begin{lstlisting}[style=terminal]
{"detail":"Authentication credentials were not provided."}
\end{lstlisting}

L'API est protégée. Chaque requête (sauf le login) doit envoyer le token.
\end{attentionbox}

\begin{retenirbox}[Le modèle Student]
Chaque étudiant a ces champs : \texttt{id}, \texttt{prenom}, \texttt{nom}, \texttt{email}, \texttt{filiere} (IABD, SIR), \texttt{niveau} (L1 à M2), \texttt{date\_naissance}, \texttt{telephone}.

C'est \textbf{différent} du modèle JSONPlaceholder (\texttt{name}, \texttt{email}, \texttt{company.name}). On va adapter le frontend.
\end{retenirbox}

% ===========================================================================
\section{Partie 2 --- Comprendre CORS}
% ===========================================================================

\begin{objectifbox}[Objectif --- 5 minutes]
Comprendre pourquoi le navigateur bloque les requêtes entre origines différentes.
\end{objectifbox}

Le frontend tourne sur \texttt{localhost:3000}. Le backend sur \texttt{localhost:8000}. Ce sont deux \textbf{origines différentes} (ports différents). Par défaut, le navigateur \textbf{interdit} à un site d'envoyer des requêtes vers une autre origine. C'est la politique CORS (\emph{Cross-Origin Resource Sharing}).

\begin{infobox}[Analogie : le gardien de l'immeuble]
Quand quelqu'un sonne, le gardien vérifie la liste des visiteurs autorisés. Si le visiteur n'est pas sur la liste, il est refusé. CORS fonctionne pareil : le backend maintient une liste d'origines autorisées.
\end{infobox}

Le backend fourni a déjà \texttt{django-cors-headers} configuré : toutes les origines sont autorisées (mode développement). Vous n'avez rien à faire côté backend. En production, on limiterait la liste.

% ===========================================================================
\section{Partie 3 --- Préparer le frontend}
% ===========================================================================

\begin{objectifbox}[Objectif --- 15 minutes]
Créer la variable d'environnement, le module API et la gestion du token.
\end{objectifbox}

Ouvrez votre projet \texttt{annuaire-next/} du Lab 5 dans VS Code. Lancez le serveur dans un \textbf{troisième terminal} :

\begin{lstlisting}[style=terminal]
cd annuaire-next
npm run dev
\end{lstlisting}

Vérifiez que l'annuaire Lab 5 fonctionne sur \texttt{http://localhost:3000}.

\subsection{Étape 1 --- Variable d'environnement}

À la \textbf{racine} du projet (pas dans \texttt{src/}), créez un fichier \texttt{.env.local} :

\begin{lstlisting}[style=terminal]
NEXT_PUBLIC_API_URL=http://localhost:8000/api
\end{lstlisting}

\begin{retenirbox}[Variables d'environnement Next.js]
\begin{center}
\begin{tabular}{lp{7cm}}
\toprule
\textbf{Préfixe} & \textbf{Accès} \\
\midrule
\texttt{NEXT\_PUBLIC\_} & Accessible côté client et serveur \\
Sans préfixe & Serveur uniquement \\
\bottomrule
\end{tabular}
\end{center}
L'URL de l'API est utilisée dans des composants \texttt{"use client"}. Elle a donc besoin du préfixe \texttt{NEXT\_PUBLIC\_}.

Au lieu d'écrire l'URL en dur dans chaque \texttt{fetch}, on écrira \texttt{process.env.NEXT\_PUBLIC\_API\_URL}. Si l'URL change (en production par exemple), on modifie un seul fichier.
\end{retenirbox}

\subsection{Étape 2 --- Module \texttt{lib/api.js}}

Créez le dossier \texttt{src/lib/} et le fichier \texttt{src/lib/api.js} :

\begin{lstlisting}[style=terminal]
mkdir -p src/lib
\end{lstlisting}

\begin{lstlisting}
// src/lib/api.js
const API = process.env.NEXT_PUBLIC_API_URL;

export const login = async (username, password) => {
  const res = await fetch(`${API}/auth/login/`, {
    method: "POST",
    headers: { "Content-Type": "application/json" },
    body: JSON.stringify({ username, password }),
  });
  if (!res.ok) {
    const data = await res.json().catch(() => ({}));
    throw new Error(
      data.error || "Identifiants invalides"
    );
  }
  return res.json();
};

export const getStudents = async (token) => {
  const res = await fetch(`${API}/students/`, {
    headers: { Authorization: `Token ${token}` },
  });
  if (!res.ok)
    throw new Error("Erreur de chargement");
  return res.json();
};

export const getStudent = async (token, id) => {
  const res = await fetch(`${API}/students/${id}/`, {
    headers: { Authorization: `Token ${token}` },
  });
  if (!res.ok)
    throw new Error("Étudiant introuvable");
  return res.json();
};

export const createStudent = async (token, data) => {
  const res = await fetch(`${API}/students/`, {
    method: "POST",
    headers: {
      "Content-Type": "application/json",
      Authorization: `Token ${token}`,
    },
    body: JSON.stringify(data),
  });
  if (!res.ok)
    throw new Error("Erreur de création");
  return res.json();
};

export const deleteStudent = async (token, id) => {
  const res = await fetch(`${API}/students/${id}/`, {
    method: "DELETE",
    headers: { Authorization: `Token ${token}` },
  });
  if (!res.ok)
    throw new Error("Erreur de suppression");
};
\end{lstlisting}

\begin{infobox}[Pourquoi centraliser les appels API ?]
Au Lab 5, les URL JSONPlaceholder étaient écrites directement dans les composants. Si l'URL change, il faut modifier chaque fichier. Avec \texttt{lib/api.js}, on change une seule ligne (\texttt{.env.local}). Chaque fonction correspond à une opération du CRUD.
\end{infobox}

\subsection{Étape 3 --- Gestion du token}

Créez \texttt{src/lib/auth.js} :

\begin{lstlisting}
// src/lib/auth.js
"use client";

export const getToken = () => {
  if (typeof window === "undefined") return null;
  return sessionStorage.getItem("token");
};

export const setToken = (token) => {
  sessionStorage.setItem("token", token);
};

export const removeToken = () => {
  sessionStorage.removeItem("token");
};

export const isLoggedIn = () => {
  return !!getToken();
};
\end{lstlisting}

Le token est stocké dans \texttt{sessionStorage}. Il disparaît quand l'onglet est fermé (plus sûr que \texttt{localStorage}).

% ===========================================================================
\section{Partie 4 --- Page de login}
% ===========================================================================

\begin{objectifbox}[Objectif --- 15 minutes]
Créer une page de connexion et protéger les pages.
\end{objectifbox}

\subsection{Créer le dossier et la page}

\begin{lstlisting}[style=terminal]
mkdir -p src/app/login
\end{lstlisting}

Créez \texttt{src/app/login/page.js} :

\begin{lstlisting}
"use client";

import { useState } from "react";
import { useRouter } from "next/navigation";
import { login } from "@/lib/api";
import { setToken } from "@/lib/auth";
import styles from "./page.module.css";

export default function LoginPage() {
  const router = useRouter();
  const [username, setUsername] = useState("");
  const [password, setPassword] = useState("");
  const [erreur, setErreur] = useState(null);
  const [enCours, setEnCours] = useState(false);

  const handleSubmit = async (e) => {
    e.preventDefault();
    setErreur(null);
    setEnCours(true);

    try {
      const data = await login(username, password);
      setToken(data.token);
      router.push("/users");
    } catch (error) {
      setErreur(error.message);
    }
    setEnCours(false);
  };

  return (
    <div>
      <h1 className={styles.titre}>Connexion</h1>
      <div className={styles.carte}>
        <form onSubmit={handleSubmit}>
          <div className={styles.groupe}>
            <label className={styles.label}
                   htmlFor="username">
              Nom d'utilisateur
            </label>
            <input className={styles.input}
              type="text" id="username"
              placeholder="prof"
              value={username}
              onChange={(e) =>
                setUsername(e.target.value)
              }
            />
          </div>
          <div className={styles.groupe}>
            <label className={styles.label}
                   htmlFor="password">
              Mot de passe
            </label>
            <input className={styles.input}
              type="password" id="password"
              placeholder="mot de passe"
              value={password}
              onChange={(e) =>
                setPassword(e.target.value)
              }
            />
          </div>
          {erreur && (
            <p className={styles.erreur}>{erreur}</p>
          )}
          <button type="submit"
                  className={styles.bouton}
                  disabled={enCours}>
            {enCours ? "Connexion..." :
                       "Se connecter"}
          </button>
        </form>
        <p className={styles.hint}>
          Identifiants de test : prof / le mot de passe affiché par setup_django.sh
        </p>
      </div>
    </div>
  );
}
\end{lstlisting}

Créez \texttt{src/app/login/page.module.css} :

\begin{lstlisting}[style=terminal]
.titre { font-size: 1.5rem; color: #2d3436;
         margin-bottom: 1.5rem; }
.carte { background: #fff;
         border: 1px solid #dfe6e9;
         border-radius: 10px; padding: 1.5rem;
         max-width: 400px; }
.groupe { margin-bottom: 1.25rem; }
.label { display: block; font-size: 0.9rem;
         font-weight: 600; color: #2d3436;
         margin-bottom: 0.4rem; }
.input { width: 100%; padding: 0.65rem 0.9rem;
         font-size: 0.95rem;
         border: 1px solid #dfe6e9;
         border-radius: 8px; }
.input:focus { outline: none;
               border-color: #0984e3; }
.erreur { color: #d63031; font-weight: 600;
          padding: 0.75rem; background: #ffeaa7;
          border-radius: 6px;
          margin-bottom: 1rem; }
.bouton { width: 100%; padding: 0.75rem;
          background: #0984e3; color: #fff;
          border: none; border-radius: 8px;
          font-size: 1rem; font-weight: 600;
          cursor: pointer; }
.bouton:hover { background: #0770c2; }
.bouton:disabled { background: #b2bec3;
                   cursor: not-allowed; }
.hint { color: #b2bec3; font-size: 0.8rem;
        text-align: center; margin-top: 1rem; }
\end{lstlisting}

\subsection{Tester le login}

Allez sur \texttt{http://localhost:3000/login}. Saisissez \texttt{prof} et le mot de passe que \texttt{setup\_django.sh} a affiché. Si la connexion réussit, vous êtes redirigé vers \texttt{/users}. Si vous tapez un mauvais mot de passe, le message d'erreur s'affiche en rouge.

\begin{attentionbox}[Redémarrer le serveur Next.js]
Après avoir créé \texttt{.env.local}, il faut \textbf{arrêter} et \textbf{relancer} \texttt{npm run dev}. Next.js ne recharge pas les variables d'environnement à chaud.
\end{attentionbox}

% ===========================================================================
\section{Partie 5 --- Modifier les pages existantes}
% ===========================================================================

\begin{objectifbox}[Objectif --- 25 minutes]
Adapter la liste, le détail, l'ajout et la Navbar pour utiliser l'API Django.
\end{objectifbox}

Les modifications suivent un principe simple : remplacer les appels directs à JSONPlaceholder par les fonctions de \texttt{lib/api.js}, et adapter les champs au modèle Student (\texttt{prenom}, \texttt{nom}, \texttt{filiere}... au lieu de \texttt{name}, \texttt{company.name}).

\subsection{Modifier \texttt{UserCard.jsx}}

Le composant reçoit maintenant un objet \texttt{student} (pas \texttt{user}). Remplacez le contenu de \texttt{src/components/UserCard.jsx} :

\begin{lstlisting}
"use client";

import Link from "next/link";
import styles from "./UserCard.module.css";

const UserCard = ({ student }) => {
  return (
    <div className={styles.card}>
      <div className={styles.header}>
        <Link href={`/users/${student.id}`}
              className={styles.lien}>
          <h3 className={styles.nom}>
            {student.prenom} {student.nom}
          </h3>
        </Link>
        <span className={styles.badge}>
          {student.filiere} {student.niveau}
        </span>
      </div>
      <p className={styles.info}>
        Email : {student.email}
      </p>
      {student.telephone && (
        <p className={styles.info}>
          Tél. : {student.telephone}
        </p>
      )}
      <Link href={`/users/${student.id}`}
            className={styles.voir}>
        Voir le profil →
      </Link>
    </div>
  );
};

export default UserCard;
\end{lstlisting}

Ajoutez la classe \texttt{.badge} dans \texttt{UserCard.module.css} :

\begin{lstlisting}[style=terminal]
.badge {
  background: #dfe6e9;
  color: #2d3436;
  padding: 0.2rem 0.6rem;
  border-radius: 12px;
  font-size: 0.75rem;
  font-weight: 600;
}
\end{lstlisting}

\subsection{Modifier la page de liste}

Remplacez le contenu de \texttt{src/app/users/page.js} :

\begin{lstlisting}
"use client";

import { useState, useEffect } from "react";
import { useRouter } from "next/navigation";
import { getStudents } from "@/lib/api";
import { getToken, isLoggedIn } from "@/lib/auth";
import UserCard from "@/components/UserCard";
import SearchBar from "@/components/SearchBar";
import styles from "./page.module.css";

export default function UserList() {
  const router = useRouter();
  const [users, setUsers] = useState([]);
  const [chargement, setChargement] = useState(true);
  const [erreur, setErreur] = useState(null);
  const [recherche, setRecherche] = useState("");
  const [triAscendant, setTriAscendant] =
    useState(true);

  useEffect(() => {
    if (!isLoggedIn()) {
      router.push("/login");
      return;
    }

    const charger = async () => {
      try {
        const data =
          await getStudents(getToken());
        setUsers(data);
      } catch (error) {
        setErreur(error.message);
      }
      setChargement(false);
    };
    charger();
  }, [router]);

  if (chargement)
    return <p>Chargement en cours...</p>;
  if (erreur)
    return <p className={styles.erreur}>
      Erreur : {erreur}</p>;

  const usersFiltres = users
    .filter((u) => {
      const terme = recherche.toLowerCase();
      const nom =
        `${u.prenom} ${u.nom}`.toLowerCase();
      return nom.includes(terme)
        || u.filiere.toLowerCase().includes(terme);
    })
    .sort((a, b) => {
      const na = `${a.nom} ${a.prenom}`;
      const nb = `${b.nom} ${b.prenom}`;
      return triAscendant
        ? na.localeCompare(nb)
        : nb.localeCompare(na);
    });

  return (
    <div>
      <h1 className={styles.titre}>
        Liste des étudiants
      </h1>
      <SearchBar recherche={recherche}
                 onRecherche={setRecherche} />
      <div className={styles.barre}>
        <span className={styles.compteur}>
          {usersFiltres.length} étudiant(s)
        </span>
        <button className={styles.boutonTri}
          onClick={() =>
            setTriAscendant(!triAscendant)
          }>
          {triAscendant ? "A → Z" : "Z → A"}
        </button>
      </div>
      {usersFiltres.map((student) => (
        <UserCard key={student.id}
                  student={student} />
      ))}
    </div>
  );
}
\end{lstlisting}

\begin{retenirbox}[Ce qui a changé par rapport au Lab 5]
\begin{enumerate}[nosep]
\item Le fetch direct est remplacé par \texttt{getStudents(getToken())} (module centralisé + token).
\item \texttt{isLoggedIn()} vérifie que l'utilisateur est connecté. Sinon, redirection vers \texttt{/login}.
\item Le filtre cherche dans \texttt{prenom + nom} et dans \texttt{filiere} (au lieu de \texttt{name}).
\item Le prop est \texttt{student} (au lieu de \texttt{user}).
\end{enumerate}
\end{retenirbox}

\subsection{Modifier la page de détail}

Remplacez le contenu de \texttt{src/app/users/[id]/page.js} :

\begin{lstlisting}
"use client";

import { useState, useEffect } from "react";
import { useParams, useRouter } from "next/navigation";
import Link from "next/link";
import { getStudent, deleteStudent } from "@/lib/api";
import { getToken, isLoggedIn } from "@/lib/auth";
import styles from "./page.module.css";

export default function UserDetail() {
  const { id } = useParams();
  const router = useRouter();
  const [student, setStudent] = useState(null);
  const [chargement, setChargement] = useState(true);
  const [erreur, setErreur] = useState(null);

  useEffect(() => {
    if (!isLoggedIn()) {
      router.push("/login");
      return;
    }

    const charger = async () => {
      try {
        const data =
          await getStudent(getToken(), id);
        setStudent(data);
      } catch (error) {
        setErreur(error.message);
      }
      setChargement(false);
    };
    charger();
  }, [id, router]);

  const handleDelete = async () => {
    if (!confirm("Supprimer cet étudiant ?"))
      return;
    try {
      await deleteStudent(getToken(), id);
      router.push("/users");
    } catch (error) {
      setErreur(error.message);
    }
  };

  if (chargement) return <p>Chargement...</p>;
  if (erreur) return (
    <div>
      <p>{erreur}</p>
      <Link href="/users">Retour</Link>
    </div>
  );

  return (
    <div>
      <Link href="/users" className={styles.retour}>
        ← Retour à la liste
      </Link>
      <div className={styles.profil}>
        <div className={styles.avatar}>
          {student.prenom.charAt(0)}
          {student.nom.charAt(0)}
        </div>
        <h1 className={styles.nom}>
          {student.prenom} {student.nom}
        </h1>
        <span className={styles.badge}>
          {student.filiere} — {student.niveau}
        </span>
      </div>
      <div className={styles.carte}>
        <h2 className={styles.section}>
          Informations
        </h2>
        <div className={styles.champ}>
          <span className={styles.label}>Email</span>
          <span>{student.email}</span>
        </div>
        <div className={styles.champ}>
          <span className={styles.label}>Tél.</span>
          <span>{student.telephone || "—"}</span>
        </div>
        <div className={styles.champ}>
          <span className={styles.label}>
            Naissance
          </span>
          <span>{student.date_naissance}</span>
        </div>
      </div>
      <button className={styles.boutonSupprimer}
              onClick={handleDelete}>
        Supprimer cet étudiant
      </button>
    </div>
  );
}
\end{lstlisting}

\begin{attentionbox}[La page de détail est devenue un composant client]
Au Lab 5, cette page était un \textbf{composant serveur} (pas de \texttt{"use client"}, fetch avec \texttt{await} directement). Maintenant, elle est un \textbf{composant client} car elle a besoin de \texttt{getToken()} (qui lit \texttt{sessionStorage}, côté client uniquement) et de \texttt{handleDelete} (événement \texttt{onClick}).

C'est un compromis : on perd le rendu serveur pour cette page, mais on gagne l'authentification et la suppression.
\end{attentionbox}

Ajoutez les classes \texttt{.badge} et \texttt{.boutonSupprimer} dans le CSS de la page de détail si elles n'y sont pas :

\begin{lstlisting}[style=terminal]
.badge {
  display: inline-block;
  background: #dfe6e9;
  padding: 0.25rem 0.75rem;
  border-radius: 20px;
  font-size: 0.85rem;
  font-weight: 500;
}
.boutonSupprimer {
  margin-top: 1rem;
  padding: 0.6rem 1.2rem;
  background: #fff;
  color: #d63031;
  border: 1px solid #d63031;
  border-radius: 8px;
  cursor: pointer;
  font-weight: 500;
}
.boutonSupprimer:hover { background: #ffeaea; }
\end{lstlisting}

\subsection{Modifier le formulaire d'ajout}

Remplacez \texttt{src/app/users/add/page.js} pour utiliser les champs du modèle Student (prénom, nom, filière, niveau, date de naissance) au lieu des champs JSONPlaceholder. Le formulaire appelle \texttt{createStudent(getToken(), form)} et vérifie \texttt{isLoggedIn()} au montage.

Les champs du formulaire :

\begin{center}
\begin{tabular}{lll}
\toprule
\textbf{Champ} & \textbf{Type} & \textbf{Obligatoire} \\
\midrule
\texttt{prenom} & Texte & Oui \\
\texttt{nom} & Texte & Oui \\
\texttt{email} & Texte (avec @) & Oui \\
\texttt{date\_naissance} & Date & Oui \\
\texttt{filiere} & Select (IABD, SIR) & Oui \\
\texttt{niveau} & Select (L1 à M2) & Oui \\
\texttt{telephone} & Texte & Non \\
\bottomrule
\end{tabular}
\end{center}

Le principe est identique au Lab 4 : état objet avec \texttt{useState}, \texttt{handleChange} avec clé dynamique, validation, \texttt{e.preventDefault()}, et redirection avec \texttt{router.push()} après succès.

\subsection{Modifier la Navbar}

Ajoutez un bouton de déconnexion dans \texttt{src/components/Navbar.jsx}. Importez \texttt{isLoggedIn} et \texttt{removeToken} depuis \texttt{@/lib/auth}. Si l'utilisateur est connecté, affichez «~Déconnexion~» (qui appelle \texttt{removeToken()} puis redirige vers \texttt{/}). Sinon, affichez un lien vers \texttt{/login}.

% ===========================================================================
\section{Partie 6 --- Tester la persistance}
% ===========================================================================

\begin{objectifbox}[Objectif --- 5 minutes]
Vérifier que le CRUD fonctionne avec de vraies données persistées.
\end{objectifbox}

\begin{enumerate}[nosep]
\item Connectez-vous avec \texttt{prof} et le mot de passe que \texttt{setup\_django.sh} a affiché.
\item La liste affiche les 10 étudiants sénégalais.
\item Allez sur \texttt{/users/add}. Ajoutez un étudiant (Ibrahima Sarr, SIR, L2).
\item Retournez sur la liste. \textbf{Le nouvel étudiant apparaît.}
\item Rechargez la page (\texttt{F5}). \textbf{Il est toujours là.} Les données sont en base.
\item Cliquez sur un étudiant, puis sur «~Supprimer~». Confirmez.
\item La liste se recharge. \textbf{L'étudiant a disparu.}
\end{enumerate}

\begin{retenirbox}[Avant / après : Lab 5 vs Lab 6]
\begin{center}
\begin{tabular}{lp{4.5cm}p{4.5cm}}
\toprule
& \textbf{Lab 5} & \textbf{Lab 6} \\
\midrule
\textbf{Source} & JSONPlaceholder (fictif) & Django REST (réel) \\
\textbf{Persistance} & Aucune & Base de données SQLite \\
\textbf{Auth} & Aucune & Token (login/logout) \\
\textbf{CRUD} & Read + faux POST & Read, Create, Delete \\
\textbf{Config URL} & En dur & Variable d'environnement \\
\textbf{Appels API} & Dispersés & Centralisés (lib/api.js) \\
\bottomrule
\end{tabular}
\end{center}
\end{retenirbox}

% ===========================================================================
\section{Aide-mémoire}
% ===========================================================================

\begin{center}
\begin{tabular}{lp{7.5cm}}
\toprule
\textbf{Concept} & \textbf{Syntaxe} \\
\midrule
Variable d'env. & \texttt{NEXT\_PUBLIC\_API\_URL=...} dans \texttt{.env.local} \\
Y accéder & \texttt{process.env.NEXT\_PUBLIC\_API\_URL} \\
\midrule
GET avec token & \texttt{headers: \{ Authorization: `Token \$\{t\}` \}} \\
POST & \texttt{method: "POST"} + headers + body \\
DELETE & \texttt{method: "DELETE"} + headers \\
\midrule
Stocker le token & \texttt{sessionStorage.setItem("token", t)} \\
Lire le token & \texttt{sessionStorage.getItem("token")} \\
Supprimer & \texttt{sessionStorage.removeItem("token")} \\
\midrule
Protéger une page & \texttt{if (!isLoggedIn()) router.push("/login")} \\
Module centralisé & \texttt{src/lib/api.js} \\
\bottomrule
\end{tabular}
\end{center}

\newpage

% ===========================================================================
\section{Résumé et conclusion}
% ===========================================================================

Dans ce lab, vous avez :
\begin{enumerate}[nosep]
\item Installé et exploré une \textbf{API Django REST} pré-remplie.
\item Compris \textbf{CORS} et pourquoi le navigateur contrôle les requêtes cross-origin.
\item Créé des \textbf{variables d'environnement} et un \textbf{module centralisé} (\texttt{lib/api.js}).
\item Implémenté l'\textbf{authentification par token} : login, stockage, envoi, déconnexion.
\item Adapté le frontend au modèle \textbf{Student} (prenom, nom, filiere, niveau...).
\item Vérifié le \textbf{CRUD complet} avec des données \textbf{persistées}.
\end{enumerate}

\begin{center}
\begin{tabular}{lcp{6.5cm}}
\toprule
\textbf{Étape} & \textbf{Statut} & \textbf{Ce que vous savez faire} \\
\midrule
\textbf{DOM et événements} & \cmark{} & Sélecteurs, événements \\
\textbf{JS moderne} & \cmark{} & Fonctions fléchées, async/await \\
\textbf{React --- les bases} & \cmark{} & Composant, JSX, props, état \\
\textbf{useEffect} & \cmark{} & Montage, dépendances \\
\textbf{Router et formulaires} & \cmark{} & Routes, Link, inputs contrôlés \\
\textbf{Next.js App Router} & \cmark{} & Routing fichier, layout, server/client \\
\textbf{API Django} & \cmark{} & CORS, auth, CRUD, lib/api.js \\
\bottomrule
\end{tabular}
\end{center}

\begin{projetbox}[Bilan du cours]
En 7 séances, vous êtes passés de la manipulation manuelle du DOM à une application Next.js complète connectée à un backend Django. Vous maîtrisez le JavaScript moderne, React (composants, état, effets, formulaires), Next.js (App Router, server/client) et la consommation d'API REST (CORS, authentification, CRUD).
\end{projetbox}

\vspace{1cm}

\begin{center}
\fbox{
\begin{minipage}{0.85\textwidth}
\centering
\vspace{0.5cm}
\textbf{\Large Lab 6 --- Consommer une API Django --- Terminé \cmark}\\[0.3cm]
\normalsize
Vous savez connecter un frontend Next.js à un backend Django,\\
gérer l'authentification et réaliser un CRUD complet.\\[0.3cm]
\textit{Du DOM manuel à la production : 7 labs, un projet.}\\[0.5cm]
\end{minipage}
}
\end{center}

\end{document}
